\section{Modelo operativo}
La operación institucional contempla gestión continua del servicio, soporte funcional por dominio y mantenimiento técnico planificado.
El ciclo operativo completo está representado en la Figura~\ref{fig:despliegue-continuidad}, donde se observan desarrollo, pruebas, producción, respaldo y recuperación.

\section{Ambientes de trabajo}
\begin{itemize}
  \item Desarrollo: configuración para iteración funcional y pruebas.
  \item Pruebas: validación controlada de cambios antes de liberación.
  \item Producción: operación estable con controles de continuidad.
\end{itemize}

\section{Despliegue y publicación}
\begin{itemize}
  \item Pipeline de publicación con control de versión documental y técnico.
  \item Parametrización por entorno y manejo seguro de credenciales.
  \item Validación post despliegue de flujos críticos de negocio.
\end{itemize}
Estos pasos corresponden directamente a la Figura~\ref{fig:despliegue-continuidad} y explican cómo se promueve una versión desde desarrollo hasta producción.

\section{Continuidad y respaldos}
\begin{itemize}
  \item Política de respaldo periódico de datos y configuraciones.
  \item Procedimientos de restauración y recuperación ante contingencias.
  \item Bitácora de incidentes y tiempos de recuperación objetivo.
\end{itemize}

\section{Monitoreo y soporte}
\begin{itemize}
  \item Seguimiento de disponibilidad, errores y rendimiento.
  \item Mesa de atencion para incidentes funcionales y tecnicos.
  \item Indicadores de nivel de servicio para mejora continua.
\end{itemize}
